% Generated draft: requires mathematical and source review.
\documentclass[11pt]{article}
\usepackage[margin=25mm]{geometry}
\usepackage{fontspec}
\setmainfont{Latin Modern Roman}
\usepackage{amsmath,amssymb,mathtools,unicode-math}
\setmathfont{Latin Modern Math}
\usepackage{enumitem,needspace,xurl,hyperref}
\setlength{\parindent}{0pt}
\setlength{\parskip}{5pt}
\setlength{\emergencystretch}{3em}
\newcommand{\N}{\mathbb{N}}
\newcommand{\Z}{\mathbb{Z}}
\newcommand{\Q}{\mathbb{Q}}
\newcommand{\R}{\mathbb{R}}
\newcommand{\C}{\mathbb{C}}
\newcommand{\F}{\mathbb{F}}
\newcommand{\A}{\mathbb{A}}
\newcommand{\PP}{\mathbb{P}}
\newcommand{\E}{\mathbb{E}}
\newcommand{\eps}{\varepsilon}
\newcommand{\dd}{\,\mathrm{d}}
\begin{document}
\section*{UnsolvedMath ID: 30000706}
\textbf{Source problem code:} catalogue entry
\subsection{Definitions and mathematical statement}
Let $f,g$ be meromorphic functions on the common complex domain intended by the source, sharing four values $a_1,a_2,a_3,a_4=\infty$. In the usual four-value convention the finite values $a_1,a_2,a_3$ are pairwise distinct, but this is not explicitly stated. Write primes for differentiation in the domain coordinate and retain all three supplied auxiliary functions:
\[
 \psi=\frac{f'g'(f-g)^2}
             {\prod_{\nu=1}^3((f-a_\nu)(g-a_\nu))},
\]
\[
 \psi_1=\frac{f'(f-g)}{\prod_{\nu=1}^3(f-a_\nu)},
 \qquad
 \psi_2=\frac{g'(f-g)}{\prod_{\nu=1}^3(g-a_\nu)}.
\]
Determine all tuples $(a_1,a_2,a_3,f,g)$ for which the quotients are meaningful meromorphic expressions and
\[
 \psi\equiv1.
\]
This is an identity of meromorphic functions, including meromorphic continuation across removable singularities, not merely equality away from a selected finite set of points. The supplied definitions give $\psi=\psi_1\psi_2$, but no additional identity for the factors is imposed.

Sharing a finite value without counting multiplicity means equality of the sets $\{z:f(z)=a_\nu\}$ and $\{z:g(z)=a_\nu\}$; sharing $\infty$ means equality of pole sets. Counting multiplicity additionally requires equality of the corresponding orders. The source does not specify which convention is intended. Nor does it identify the domain, for example the whole complex plane versus a restricted domain, or explicitly state the usual nonconstancy and distinct-value assumptions. A function identically equal to a denominator value must be excluded or separately interpreted rather than treating an identically undefined quotient as an identity.

The classification target is all tuples with these shared-value conditions and the exact normalization $\psi=1$. No equivalence under changes of values or domain variable, no growth hypothesis, and no counting-multiplicities strengthening is silently added. The unspecified domain and sharing conventions must be recovered for a complete classification.
\subsection{Short English statement}
Classify the value triples and meromorphic pairs sharing those values and infinity for which the displayed Mues auxiliary function is identically one. All three auxiliary definitions are preserved, with domain and multiplicity-sharing conventions left explicit for source review.
\subsection{Sources}
\begin{itemize}
\item Local UnsolvedMath registry, permanent ID 30000706, code catalogue entry, ``Four-Value-Sharing Meromorphic Functions with $\psi=1$.'' Proposed year 2007; proposer unspecified. The assessment identifies this as 2007 Problem 1 and names the Mues auxiliary function.
\item \emph{Normal Families and Complex Dynamics} (2007), Oberwolfach Reports, pp.\ 487--548, DOI 10.4171/owr/2007/09; \url{https://doi.org/10.4171/owr/2007/09}.
\item Supplied ``later partial four-value results'': \url{https://doi.org/10.1016/S0252-9602(17)30234-5}.
\item Statement status corrected\_verified; the supplied note says the definitions and classification target are explicit and Unicode primes, subscripts, and products were normalized.
\item Literature triage checked 2026-08-21: open. It says later literature gives partial cases but no complete classification of this four-value-sharing normalized equation. Exact hypotheses and current status require independent verification. Created 2026-08-25; updated 2026-09-16.
\item Registry curation by Ulam AI / UnsolvedMath Contributors, CC BY 4.0; underlying report and literature provenance remain as supplied. Unreviewed statement-only draft.
\end{itemize}
\end{document}
